% Part: computability
% Chapter: recursive-functions
% Section: notations-pr-functions

\documentclass[../../../include/open-logic-section]{subfiles}

\begin{document}

\olfileid{cmp}{rec}{not}
\olsection{Notasi Rekursi Primitif}

Salah siji paedahe gambaran induktif kang cetha ngenani fungsi
rekursif primitif yaiku kita bisa nggambarake fungsi kasebut
kanthi sistematis. Umpamane, kita bisa menehi “notasi” kanggo
saben fungsi kasebut kaya mangkene. Gunakna lambang $\Zero$,
$\Succ$, lan $\Proj{n}{i}$ kanggo fungsi nol, penerus, lan
proyeksi. Saiki upamana $h$ ditegesi kanthi komposisi saka
fungsi $k$ panggonan~$f$ lan fungsi $n$ panggonan $g_0$,
\dots,~$g_{k-1}$, lan kita wis menehi notasi $F$, $G_0$,
\dots,~$G_{k-1}$ marang fungsi-fungsi kasebut. Banjur,
nganggo lambang anyar $\fn{Comp}_{k,n}$, kita bisa nulis
fungsi $h$ minangka $\fn{Comp}_{k,n}[F,G_0,\dots,G_{k-1}]$.

Kanggo fungsi kang ditegesi kanthi rekursi primitif,
kita bisa nganggo notasi kang padha. Upamana fungsi
$(k+1)$ panggonan~$h$ ditegesi kanthi rekursi primitif
saka fungsi $k$ panggonan~$f$ lan fungsi $(k+2)$ panggonan~$g$,
lan notasi kanggo $f$ lan~$g$ yaiku $F$ lan~$G$
kanthi urutan kuwi. Banjur notasi kanggo~$h$ yaiku
$\fn{Rec}_k[F,G]$.

Elinga yen fungsi panjumlahan ditegesi kanthi rekursi primitif
kaya mangkene:
\begin{align*}
  \Add(x_0, 0) & = \Proj{1}{0}(x_0) = x_0\\
  \Add(x_0, y+1) & = \Succ(\Proj{3}{2}(x_0, y, \Add(x_0, y))) = \Add(x_0, y) +1
\end{align*}
Ing kene peran~$f$ dimainake dening $\Proj{1}{0}$, lan peran~$g$
dimainake dening $\Succ(\Proj{3}{2}(x_0, y, z))$, kang diwenehi
notasi $\fn{Comp}_{1,3}[\Succ,\Proj{3}{2}]$ amarga fungsi iki
asile komposisi saka fungsi $1$ panggonan~$\Succ$
lan fungsi $3$ panggonan~$\Proj{3}{2}$. Kanthi tatanan iki,
kita bisa nulis fungsi panjumlahan minangka
\[
\fn{Rec}_1[\Proj{1}{0},\fn{Comp}_{1,3}[\Succ,\Proj{3}{2}]].
\]
Notasi kasebut kadhangkala migunani, umpamane nalika
ngenumerasi fungsi rekursif primitif.

\begin{prob}
  Wenehana notasi rekursif primitif kang lengkap kanggo~$\Mult$.
\end{prob}

\end{document}
